\documentclass[12pt,a4paper]{article}
\usepackage{fontspec}
\usepackage{polyglossia}
\setmainfont{TeX Gyre Termes}
\newfontfamily\cyrillicfont{TeX Gyre Termes}
\setdefaultlanguage{ukrainian}
\setotherlanguage{english}
\usepackage{amsmath,amssymb,geometry,graphicx}
\geometry{margin=25mm}
\setlength{\parindent}{1.25cm}
\setlength{\parskip}{0.25em}
\begin{document}
\title{2025 EM Parallel implementation of polypoint transformations with adjacent triangle plane intersections}
\maketitle
У попередньому дослідженні було проаналізовано 3 методи задання топології обєкта перетворення методом ПТП: метод перетину площин трикутників та його нормалей, метод ортогональних площин для кожної вершини та метод задання через перетин площин дотичних трикутників. Було зроблено висновок, що найбільш преспективним зважаючи на швидкодію та ефективність деформації виявився метод задання геометрії через перетин площин дотичних трикутників. Дане дослідження присвячено тому, щоб покращити ефективність ПТП з використанням цього методу.
\par
Сутність методу перетинаючихся площин полягає в тому, що кожен вузол тріангульованого каркасу задається через список площин, утворених з трикутників, вершиною кожного з яких є вузол (Рис.
\par
Вершин А, задана площинами, утвореними з
\par
Після застосування ПТП для кожної площини із списку трикутників отримаємо список з перетворених площин. Розглянемо найпростіший випадок, коли вершина задається через три площини, наприклпд
\par
перетворені площини на основі яких необхідно відшукати перетворену вершину А
\par
після застосування ПТП та шукана вершина А
\par
У дослідженні Сидоренко Ю.В. та Бадаєва Ю
\par
на 99\% певен, що не було такого у Юрія Івановича. Чисельну оптимізацію запропонував я минулого року, щоб мати хоч якийсь результат. Це взагалі не мало б бути опубліковано) пропонувалось знаходити вершину А
\par
методом оптимізації, тобто шукаючи точку, яка б була мінімально віддалена від кожної із площин.
\par
ФОРМУЛА, яка показує як записується мінімум.
\par
Я не знаю, як правильно записати
\par
Дійсно, такий підхід дозволяє знайти таку точку незалежно від кількості площин, але недоліком є його швидкодія.
\par
Іншим способом знаходження вершини є метод найменших квадратів.
\par
Я не знаю, як правильно записати
\par
Ми пропонуємо використати метод псевдоінверсної матриці Мура Пенроуза
\par
для знаходження найближчої точки до всіх перетворених площин, якими задана вершина A
\par
псевдообернена матриця Мура-Пенроуза до
\par
Ми зробили заміри часу розрахунків деформації тріангульованої моделі тора методом ПТП з різною кількістю полігонів (рис.
\par
3. Графік залежності методу пошуку вершина A
\par
від часу виконання в секундах
\par
З графіку видно, що на малій кількості трикутників час виконання приблизно однаковий, однак із збільшенням кількості обчислень метод мінімізації працює значно повільніше за pinv та lstsq. Метод pinv показує найкращі результати на різній кількості трикутників, хоча має невелику перевагу відносно lstsq. Це одна і та сама задача, оця матриця Мура-Пенроуза – це спосіб розв’язку задачі найменшої суми квадратичних відхилень. Тому порівнювати їх якось дивно. Який спосіб розв’язку ти називаєш lstsq?
\par
Далі я думав додати ще паралелізацію, де буде розбиття на чанки і запуск обчислень для кожного чанку в окремому процесі (бо це Python). Висновком буде, що алгоритм ПТП добре паралелиться і буде видно приріст ефективності...
\par
Так в цьому і ідея. Писати, що аналітичний спосіб розв’язку квадратичної задачі швидше за чисельний – не треба. Це ніхто ніколи і не оскаржував.
\par
Забий на Пітон, куди міг він тебе довів. Якщо знаєш С++ з OMP чи TBB краще за пітон із numba і numpy, переходь на С++. Немає жодного сенсу дивитися на швидкодію екпериментального коду.
\par
\end{document}
